\section{Gestion hebdomadaire et suivi de l'expérimentation culturale}
La réussite de l’expérimentation repose sur un suivi régulier des cultures tout au long de leur cycle de développement. Ce suivi a porté sur l’irrigation, l’entretien des cultures, le suivi phytosanitaire et l’évolution des différents plots.

\subsection{Protocole de suivi des cultures}
Dès le lendemain du semis, le suivi des cultures a été effectué trois fois par jour (7~h–8~h, vers 13~h et 18~h–19~h). Pendant la première semaine, une irrigation régulière a été assurée jusqu’à la levée et à l’installation des jeunes plants. Les observations réalisées ont été renseignées dans l’agenda de culture et accompagnées de photographies.

\begin{table}[H]
    \centering
    \small
    \caption{Protocole de suivi des cultures durant l'expérimentation}
    \label{tab:protocole_suivi}
    \begin{tabular}{|p{3.2cm}|p{5.4cm}|p{5.4cm}|}
        \hline
        \textbf{Éléments suivis} & \textbf{Opérations réalisées} & \textbf{Données recueillies} \\
        \hline
        Irrigation et suivi hydrique & Arrosage régulier, observation du comportement du sol et adaptation des irrigations & Apports d'eau, temps d'infiltration, pluies et comportement du sol \\
        \hline
        Suivi phénologique & Observation de la germination, de la levée et de l'évolution des feuilles & Nombre de poquets levés, apparition des feuilles, déformations et changements de couleur \\
        \hline
        Entretien des cultures & Observation de l'état général des plants & Anomalies et mortalités observées \\
        \hline
        Suivi phytosanitaire & Observation des ravageurs, des dégâts et des symptômes de maladies ; application des traitements nécessaires & Ravageurs observés, symptômes, dégâts et traitements appliqués \\
        \hline
        Documentation & Enregistrement des observations dans l'agenda de culture et prise de photographies & Historique de l'évolution des cultures et des interventions \\
        \hline
    \end{tabular}
\end{table}

\subsection{Entretien des cultures}
L’entretien des cultures a été effectué régulièrement sur l’ensemble des six plots en fonction de leur état et du stade de développement des cultures. Il comprenait principalement \textbf{le désherbage, le binage, le sarclage, le sarclo-binage, la délimitation de cuvettes, le buttage, le tuteurage, l’irrigation, la fertilisation et la récolte} afin de maintenir des conditions de développement favorables aux cultures.

Le \textbf{désherbage et le sarclage} ont consisté à éliminer les adventices afin de limiter leur concurrence pour la lumière, l’eau, l’espace et les nutriments. Le \textbf{binage}, souvent associé au sarclage pour former le sarclo-binage, a permis d’ameublir la couche superficielle du sol, d’améliorer son aération et de faciliter l’infiltration de l’eau tout en réduisant son évaporation. La \textbf{délimitation des cuvettes} a favorisé la rétention de l’eau dans la zone racinaire et limité le ruissellement. Le \textbf{tuteurage} du gombo a servi de soutien aux plants pour réduire les risques de verse et de casse, tandis que le \textbf{buttage} du navet a contribué à protéger son tubercule et à favoriser son développement.
Certaines opérations d’entretien ont été adaptées aux stades de développement et aux exigences spécifiques de chaque culture. L’habillage du GOM TS a permis de réduire l’excès de feuillage et de favoriser la circulation de la lumière et de l’air. Le buttage du navet a été effectué lors de la formation du tubercule. Enfin, la récolte progressive a été réalisée selon le stade de maturité propre au gombo, au navet et à la laitue, permettant de suivre leur évolution et leur rendement.
\begin{figure}[H]
    \centering
    \imgstd{images/29_binage_sarclo_tuteurage.jpg}
    \caption{Actions respectives de Sarclo-binage, de Binage et de tuteurage. }
    \label{fig:entretient3}
\end{figure}

\subsection{Apport journalier d’eau}
Le pilotage de l’irrigation repose sur le suivi de l’humidité du sol et sur l’estimation des besoins en eau des cultures à partir de l’ETc calculée dans le classeur Excel.

\begin{table}[H]
    \centering
    \small
    \caption{Protocole de suivi et de pilotage de l'irrigation}
    \label{tab:pilotage_irrigation}
    \begin{tabular}{|p{4.2cm}|p{9.8cm}|}
        \hline
        \textbf{Éléments suivis} & \textbf{Modalités appliquées} \\
        \hline
        Plots concernés & GOM TS, NAV TS et LAI TS \\
        \hline
        Points de mesure & Deux points par plot \\
        \hline
        Outil de mesure & Capteur 8-en-1 \\
        \hline
        Fréquence de mesure & Deux fois par jour \\
        \hline
        Critère de déclenchement & Irrigation lorsque $S_a < S_c$ \\
        \hline
        Absence d'irrigation & Lorsque $S_a \geq S_c$ \\
        \hline
        Calcul de la dose & $ET_c$ calculée quotidiennement dans le classeur Excel à partir des données météorologiques collectées ou disponibles en prévision \\
        \hline
        Répartition de la dose & Deux apports : matin et soir \\
        \hline
        Durée d'application & Jusqu'à la fin de l'expérimentation \\
        \hline
    \end{tabular}
\end{table}

\begin{figure}[H]
    \centering
    \imgstd{images/30_collecte_hebdomadaire.jpg}
    \caption{Collecte des données sur le plot LAI TS. }
    \label{fig:collecte3}
\end{figure}

\subsection{Stratégie de lutte contre les maladies et ravageurs}
Au cours de notre expérimentation, divers ravageurs et maladies ont été observés sur le gombo, le navet et la laitue. Pour chaque cas, les symptômes ou dégâts, les agents responsables et les méthodes de lutte appliquées ont été identifiés et consignés.

\subsubsection{Ravageurs observés et méthodes de lutte}
Au cours de l’expérimentation, plusieurs ravageurs ont été observés sur les cultures, notamment le gombo et le navet. Les observations réalisées et les méthodes de lutte appliquées sont présentées dans le tableau suivant.

{\small
\begin{longtable}{|p{2.6cm}|p{3.4cm}|p{4.6cm}|p{4cm}|}
    \caption{Ravageurs observés, dégâts et méthodes de lutte}
    \label{tab:ravageurs_lutte} \\
    \hline
    \textbf{Ravageurs observés} & \textbf{Symptômes et dégâts constatés} & \textbf{Méthodes de lutte appliquées} & \textbf{Résultats observés} \\
    \hline
    \endfirsthead
    \hline
    \textbf{Ravageurs observés} & \textbf{Symptômes et dégâts constatés} & \textbf{Méthodes de lutte appliquées} & \textbf{Résultats observés} \\
    \hline
    \endhead
    Cochenilles farineuses (gombo) & Amas blancs et cotonneux sous les feuilles, jaunissement foliaire et présence de miellat. & Retrait manuel des cochenilles et utilisation ponctuelle d'un jet d'eau. & Les cochenilles sont éliminées des feuilles après l'intervention. Cependant, leur présence est de nouveau constatée quelques jours plus tard. \\
    \hline
    Cicadelle écumeuse (gombo) & Présence de mousse blanche sur les feuilles, protégeant les larves. & Jet d'eau pour éliminer la mousse et retrait manuel à l'éponge des larves et de la mousse. & Les larves sont délogées et éliminées lors du traitement. Toutefois, de nouvelles présences sont observées quelques jours après l'intervention. \\
    \hline
    Chenilles défoliatrices (gombo et navet) & Trous dans les feuilles et défoliation. & Retrait manuel et pulvérisation de solutions à base de neem ou d'ail, généralement en fin de journée. La solution de neem était préparée avec 500~g de feuilles pilées dans 10~L d'eau et une petite quantité de savon. La solution d'ail était préparée avec 500~g de gousses dans 9~L d'eau, complétée par une solution savonneuse. & L'odeur des solutions semble éloigner temporairement les chenilles, mais celles-ci réapparaissent après quelques jours. Lors des applications répétées d'une même solution, son effet répulsif semble diminuer progressivement. \\
    \hline
    Criquets (gombo et navet) & Consommation et sectionnement des tissus foliaires. & Pulvérisation d'une solution à base d'ail sur le feuillage. & La solution semble efficace au début, en éloignant les criquets. Toutefois, son effet répulsif diminue avec le temps, et les criquets finissent par revenir. \\
    \hline
\end{longtable}
}

Les ravageurs ont été identifiés à partir des observations effectuées sur les cultures. Les méthodes de lutte ont été adaptées aux types de dégâts constatés. Les photographies suivantes illustrent les ravageurs observés durant l’expérimentation.

    \begin{figure}[H]
    \centering
    \imgstd{images/31_cochenille_farineuse.jpg}
    \caption{Présence des cochenilles farineuses sur une feuilles de gombo. }
    \label{fig:cochenille3}
\end{figure}

    \begin{figure}[H]
    \centering
    \imgstd{images/32_cicadelle_ecumeuse.jpg}
    \caption{Cicadelle écumeuse sur une feuille de gombo : (a) mousse protectrice ; (b) larve extraite de la mousse. }
    \label{fig:cicadelle}
\end{figure}

\begin{figure}[H]
    \centering
    \imgstd{images/33_chenille.jpg}
    \caption{Ravageurs observés : 2 types de chenilles (en haut) et 2 types de criquets (en bas). }
    \label{fig:chenille3}
\end{figure}

\subsubsection{Maladies observées et méthodes de lutte}
L'observation a montré des symptômes compatibles avec ceux des maladies fongiques.

\begin{itemize}[label=$\bullet$]
   	\item \textbf{Les maladies fongiques : } sont causées par le développement de champignons pathogènes, accéléré par la combinaison d'une intense chaleur et d'une forte humidité. Trois types de maladies fongiques ont particulièrement été observés au cours de notre expérimentation : \textbf{la fusariose, le phytophthora et l'anthracnose.}

    \begin{itemize}[label=$\circ$]
   	\item \textbf{La fusariose}
ou flétrissement vasculaire, est principalement causée par un champignon appelé \textbf{Fusarium oxysporum}. Il pénètre l'organisme par les racines et colonise les vaisseaux de conduction de la sève (le xylème), sécrétant des toxines qui obstruent progressivement les canaux. Elle se manifeste par une chlorose et un flétrissement soudain, qui commencent généralement d'un côté de la plante ; une section de la tige révèle des vaisseaux internes brunis ou noircis. La sève ne circulant plus, la plante se dessèche en interne et meurt rapidement. Cette maladie est favorisée par un sol réchauffé (25 à 32°C), un pH acide et est encouragée par un stress hydrique causé par des périodes sèches alternées avec des inondations. De plus, le champignon a la capacité de survivre plusieurs années dans le sol, sous forme de spores résistantes.

    \item Le \textbf{phytophthora} responsable de la pourriture du collet et des racines, n'est pas un champignon à proprement parler mais un oomycète, un pathogène aquatique qui exploite un sol gorgé d'eau pour infiltrer les racines avec ses spores en mouvement. Il détruit les cellules racinaires afin d'absorber les matières organiques. Les signes de la maladie se manifestent par un brunissement et un ramollissement du collet, des racines noircies et visqueuses, et un flétrissement soudain des feuilles supérieures, entraînant l'effondrement et la mort de la plante en quelques jours. Cette maladie est favorisée par un excès d'eau (sols compacts, argileux ou trop irrigués), des températures tempérées à élevées  (15 à 25°C) et le manque d'oxygène au niveau des racines.

    \item \textbf{L'anthracnose} est une maladie foliaire et fruitière causée principalement par des champignons du genre \textbf{Colletotrichum}. Elle se maintient sur les résidus végétaux ou les graines et se répand par dispersion des spores lors des pluies. Elle se manifeste par des taches circulaires concaves, allant du brun au noir, parfois entourées d'un halo jaune sur les feuilles, ainsi que par des lésions noires et sèches sur les fruits (les capsules de gombo), qui les rendent non comestibles et non vendables. Une attaque sévère peut provoquer un dessèchement total et une chute des feuilles, engendrant une baisse significative du rendement. Elle est favorisée par une forte humidité (une pluie fréquente, une rosée matinale étendue, ou une irrigation par aspersion), une chaleur oppressante (20 à 28°C), ainsi qu'une densité de plantation excessive limitant la circulation de l'air.
    
    \item Méthode de lutte : Les plantes déjà atteintes n’ont pas pu être sauvées. Pour les plantes encore saines, une lutte préventive a été réalisée à l'aide de la solution d’ail et d’un traitement au bicarbonate, préparés à raison d’une cuillère à café par litre d’eau. Les traitements sont appliqués sur les feuilles en fin de journée (19h--20h).
	\end{itemize}
\end{itemize}

    \begin{figure}[H]
    \centering
    \imgstd{images/34_fongique.jpg}
    \caption{Gombo atteint d'une maladie fongique montrant un collet et des racines pourries. }
    \label{fig:fongique3}
\end{figure}

\subsection{Évolutions des différents plots}
Tout au long du cycle expérimental, les cultures ont fait l’objet d’une surveillance visuelle et qualitative quotidienne sur l'ensemble des plots, portant notamment sur \textbf{la vigueur, le feuillage et la croissance} des cultures. 
Ces observations ont été complétées par des photographies, des mesures régulières et le suivi des stades phénologiques clés, permettant une comparaison entre les plots T et TS. Les attaques de ravageurs et les maladies fongiques ont également été surveillées et consignées dans notre carnet de suivi expérimental.


















